\section{\BGA{} arithmetized in \QGA}
\label{sec:bga}

\texttt{BGA\_on\_GA.thy} encodes basic grounded arithmetic as \QGA{} natural
numbers and proves its proof checker sound with respect to the encoded
semantics.  The file is a tower of locales, so the whole argument is carried
out under a fixed, small set of assumptions about the encoding and no concrete
G\"odel numbering appears in it.

This section walks through the locales in order and explains every constant
they fix.  Throughout, $\Gamma \seq \varphi$ is a \BGA{} judgment, $\pi$ and
$\rho$ range over encoded proofs, $A$ over assignments, and $D$ over the
background definition list.  Notation like $q[v_i \mapsto a]$ abbreviates
\isa{subst\_F q i a}.

\subsection{The system being encoded}

\BGA{} has terms and formulas, and no logical connectives at all.

\[
\begin{array}{lcl}
t &::=& v_j \mid 0 \mid \Suc{t} \mid \Pre{t} \mid t\,0?\,t : t \mid d_i(t, t)\\
\varphi &::=& t = t \mid t \neq t
\end{array}
\]

\noindent
A variable $v_j$ is a natural number $j$, and its value under an assignment is
the $j$th entry of the assignment list.  \BGA{} has no binders, so nothing
scopes over a variable and the index is simply a position.  Conditional
evaluation is if-zero-else, which keeps terms and formulas from being mutually
recursive.  Function application is restricted to two arguments, and the
functions come from a fixed background list $D$, where position $i$ holds the
body of $d_i(v_0, v_1)$.  Disequality is a primitive, so $a \neq b$ is not the
negation of $a = b$; \BGA{} has no negation to form.  A judgment is a sequent
$\Gamma \seq \varphi$ with $\Gamma$ a finite list of formulas.

\subsection{Codes}

There is one \QGA{} type, \isa{num}.  The type synonyms below are
documentation.

\thy{bga-types}

\noindent
A judgment is the pair \isa{⟨G, c⟩}, written \isa{G ⊩ c}, with \isa{hyp\_of}
and \isa{conc\_of} as projections.  A hypothesis set is a list of formulas and
a proof is a list of judgments, read head first, each justified by the
judgments after it.

\subsection{The locale tower}
\label{sec:tower}

\begin{center}
\begin{tabular}{@{}lll@{}}
\toprule
\textbf{Locale} & \textbf{Extends} & \textbf{Fixes}\\
\midrule
\isa{suff\_syntax}    & ---                   & \isa{mk\_eq}, \isa{mk\_neq}, \isa{dfns}, \isa{is\_valid\_proof}\\
\isa{suff\_semantics} & \isa{suff\_syntax}    & \isa{evals}, \isa{sat\_fm}, \isa{sat\_hyp}\\
\isa{consistent}      & \isa{suff\_semantics} & ---\\
\midrule
\isa{bga\_encoding}   & ---                   & the six encoder operations, the three\\
                      &                       & freshness predicates\\
\isa{bga\_dfns}       & \isa{bga\_encoding}   & \isa{dfns}, \isa{dfn\_is}\\
\isa{bga\_fuel\_semantics} & \isa{bga\_dfns}  & \isa{eval\_fuel}\\
\isa{bga\_subst}      & \isa{bga\_encoding}   & \isa{subst\_T}, \isa{subst\_F}, \isa{subst\_body}\\
\isa{bga\_fuel\_subst\_semantics} & \isa{bga\_fuel\_semantics}, & \isa{asn\_put}\\
                      & \isa{bga\_subst}      & \\
\midrule
\isa{bga\_subst\_rule}  & \isa{bga\_subst}    & \isa{check\_template}, \isa{rep\_vars\_T},\\
                        &                     & \isa{rep\_vars\_F}, \isa{find\_phi}, \isa{find\_eq},\\
                        &                     & \isa{check\_subst}\\
\isa{bga\_eq\_rule}     & \isa{bga\_encoding} & \isa{check\_eq\_rules}, \isa{check\_neq\_rules}\\
\isa{bga\_ind\_rule}    & \isa{bga\_subst\_rule} & \isa{check\_ind\_template},\\
                        &                     & \isa{find\_ind\_base}, \isa{check\_ind}\\
\isa{bga\_struct\_rule} & \isa{bga\_encoding} & \isa{find\_cut}, \isa{check\_cut},\\
                        &                     & \isa{find\_struct}, \isa{check\_struct}\\
\isa{bga\_app\_rule}    & \isa{bga\_subst\_rule}, & \isa{app\_try}, \isa{app\_y}, \isa{app\_x},\\
                        & \isa{bga\_dfns}     & \isa{app\_d}, \isa{check\_app}\\
\isa{bga\_proof\_check} & the five rule locales & \isa{valid\_step}, \isa{check\_list},\\
                        &                     & \isa{is\_valid\_proof}\\
\midrule
\isa{bga\_fuller}       & \isa{bga\_fuel\_subst\_semantics}, & ---\\
                        & \isa{bga\_proof\_check} & \\
\bottomrule
\end{tabular}
\end{center}

\noindent
Two locales are defined in the file and not used by \isa{bga\_fuller}:
\isa{bga\_semantics}, which fixes the exact evaluator \isa{eval}, and
\isa{bga\_subst\_semantics}, which combines it with substitution.  They are
described in \S\ref{sec:fuel} because the exact evaluator is the clearest
statement of what the encoded semantics means.

Almost every constant is fixed with an assumption of the form \isa{f x := body}
rather than being defined.  A locale cannot make a recursive definition, and in
\QGA{} it does not need to: fixing a constant together with its \isa{:=}
equation is exactly what \S\ref{sec:defmech} admits, and the equation is all
the proofs ever use.

\subsection{The consistency skeleton}
\label{sec:skeleton}

Before any encoding is fixed, the shape of the final argument is stated.

\thy{bga-suff-syntax}

\begin{center}
\begin{tabular}{@{}ll@{}}
\toprule
\isa{mk\_eq :: tm ⇒ tm ⇒ fm} & builds the code of $a = b$ from the codes of $a$ and $b$\\
\isa{mk\_neq :: tm ⇒ tm ⇒ fm} & the same for $a \neq b$\\
\isa{dfns :: dfn} & the code of the background definition list $D$\\
\isa{is\_valid\_proof :: pf ⇒ jdg ⇒ o} & ``$\pi$ is a proof of $J$''\\
\bottomrule
\end{tabular}
\end{center}

\noindent
The assumptions are only habeas quid: encoded formulas are numbers, and
\isa{is\_valid\_proof p J} is a grounded formula whenever \isa{p} and \isa{J}
are numbers.  The second is the grounded-logic form of decidability of proof
checking, and it is what makes the concluding proof by contradiction available.

\thy{bga-suff-semantics}

\begin{center}
\begin{tabular}{@{}ll@{}}
\toprule
\isa{evals :: tm ⇒ asn ⇒ val ⇒ o} & ``$t$ evaluates to $r$ under $A$''\\
\isa{sat\_fm :: fm ⇒ asn ⇒ o} & ``the formula coded by $f$ holds under $A$''\\
\isa{sat\_hyp :: hyp ⇒ asn ⇒ o} & ``every formula in $\Gamma$ holds under $A$''\\
\bottomrule
\end{tabular}
\end{center}

\noindent
Three assumptions: the empty hypothesis list is satisfied by every assignment,
evaluation is deterministic, and a satisfied encoded equation yields a common
value for its two sides while a satisfied encoded disequality yields two
distinct values.

\thy{bga-consistent}

\noindent
\isa{consistent} adds soundness of the checker and derives syntactic
consistency.  The derivation is short.  Assume both proofs exist.  The
statement is grounded by \isa{proof\_bool}, so \isa{contradiction} applies.
Soundness under the empty hypothesis list and the assignment \isa{0} gives
\isa{sat\_fm (mk\_eq a b) 0} and \isa{sat\_fm (mk\_neq a b) 0}.  The first
produces a value \isa{q} reached by both \isa{a} and \isa{b}, the second
produces \isa{x} and \isa{y} with \isa{x ≠ y}, determinism forces
\isa{x = q = y}, and \isa{exF} closes the goal.  The habeas quid premises
\isa{a N}, \isa{b N}, \isa{p1 N} and \isa{p2 N} are needed before the formula
being contradicted can be shown grounded.

\subsection{Encoded syntax}

Terms and formulas are tagged pairs, with six term tags and two formula tags.

\thy{bga-tags}

\begin{center}
\begin{tabular}{@{}llll@{}}
\toprule
Tag & Constructor & Notation & Payload\\
\midrule
\isa{T\_VAR}  & variable      & $v_j$              & $j$\\
\isa{T\_ZERO} & zero          & $0$                & ignored\\
\isa{T\_SUC}  & successor     & $\Suc{a}$          & $a$\\
\isa{T\_PRED} & predecessor   & $\Pre{a}$          & $a$\\
\isa{T\_IFZ}  & if-zero-else  & $c\;0?\,a:b$       & $\langle c, \langle a, b\rangle\rangle$\\
\isa{T\_APP}  & application   & $d_i(a,b)$         & $\langle i, \langle a, b\rangle\rangle$\\
\midrule
\isa{F\_EQ}   & equality      & $a = b$            & $\langle a, b\rangle$\\
\isa{F\_NEQ}  & disequality   & $a \neq b$         & $\langle a, b\rangle$\\
\bottomrule
\end{tabular}
\end{center}

\subsection{\isa{bga\_encoding}}
\label{sec:encoder-interface}

\thy{bga-enc-fixes}

\begin{center}
\begin{tabular}{@{}ll@{}}
\toprule
\isa{tag\_T :: tm ⇒ tmtag} & the constructor tag of a term code\\
\isa{load\_T :: tm ⇒ tm}   & the payload of a term code\\
\isa{pack\_T :: tmtag ⇒ tm ⇒ tm} & builds a term code from tag and payload\\
\isa{tag\_F}, \isa{load\_F}, \isa{pack\_F} & the same three for formula codes\\
\isa{fresh\_T :: num ⇒ tm ⇒ o} & every variable in the term has index below $k$\\
\isa{fresh\_F :: num ⇒ fm ⇒ o} & the same for a formula\\
\isa{fresh\_H :: num ⇒ hyp ⇒ o} & the same for every formula in a list\\
\bottomrule
\end{tabular}
\end{center}

\noindent
Freshness is defined by recursion over the tag, and in mathematical notation it
reads
\[
\fn{fresh}_T(k, t) =
\begin{cases}
j < k                                        & t = v_j\\
\top                                         & t = 0\\
\fn{fresh}_T(k, a)                           & t = \Suc{a} \text{ or } t = \Pre{a}\\
\fn{fresh}_T(k, c) \wedge \fn{fresh}_T(k, a) \wedge \fn{fresh}_T(k, b)
                                             & t = (c\,0?\,a:b)\\
\fn{fresh}_T(k, a) \wedge \fn{fresh}_T(k, b) & t = d_i(a, b)\\
\bot                                         & \text{tag out of range}
\end{cases}
\]
with $\fn{fresh}_F(k, \varphi)$ the conjunction over the two sides of
$\varphi$ and $\fn{fresh}_H(k, \Gamma)$ the conjunction over the members of
$\Gamma$.  The last line is the treatment of a number whose tag is not one of
the six: it is fresh for no $k$.  Freshness is used once, in the induction rule
of \S\ref{sec:ind-rule}, to guarantee that the induction variable does not
occur in the hypotheses.

The assumptions on the interface are the following, and they are the only
properties of the encoding used anywhere in the development.

\thy{bga-enc-axioms}

\begin{center}
\begin{tabular}{@{}ll@{}}
\toprule
habeas quid & \isa{pack}, \isa{tag} and \isa{load} preserve \isa{N}\\
injectivity & \isa{tag (pack t L) = t} and \isa{load (pack t L) = L}\\
retraction & \isa{pack (tag f) (load f) = f}\\
decrease & \isa{f ≠ 0 ⟹ load f < f = 1}, payloads are strictly smaller\\
base case & \isa{tag\_T 0 = T\_ZERO}, the code $0$ is the term $0$\\
monotonicity & \isa{pack} is monotone in its payload\\
\bottomrule
\end{tabular}
\end{center}

\begin{remarkx}[Injective, not surjective]
What the concrete encoding of \S\ref{sec:encode} provides, and all that is
used, is an \emph{injective} map from \BGA{} syntax into $\mathbb{N}$.  It is
not surjective: a number whose tag lies outside the constructor range encodes
no \BGA{} term, and no attempt is made to exclude such numbers or to recognize
the well-formed ones.  The assumption that does hold for every number is the
retraction property \isa{pack (tag f) (load f) = f}, which follows from the
Cantor pairing being a bijection on $\mathbb{N}$ and says only that a code can
be taken apart and put back together.  Ill-formed codes are harmless: the
evaluator gives them the application semantics, freshness rejects them, and
soundness never needs them to denote anything.
\end{remarkx}

\noindent
The decrease and monotonicity assumptions are what allow
\isa{strong\_induction} to serve as structural induction over encoded terms.
Decrease gives the descent: the payload of a nonzero code is a strictly smaller
number, so recursion into subterms is bounded numerical recursion.
Monotonicity is used differently, to justify the bounds of the searches in
\S\ref{sec:subst-rule}.  The first lemma proved from them is
\isa{fresh\_T\_bool}, that freshness is a grounded formula.

\subsection{\isa{bga\_dfns}}

\thy{bga-dfns}

\begin{center}
\begin{tabular}{@{}ll@{}}
\toprule
\isa{dfns :: dfn} & the code of the definition list $D$, a list of term codes\\
\isa{dfn\_is :: dfn ⇒ num ⇒ tm ⇒ o} & ``$d$ is in range, $D_d$ is $b$, and $b$ has arity $\leq k$''\\
\bottomrule
\end{tabular}
\end{center}

\noindent
Written out, $\fn{dfn\_is}(d, k, b)$ holds when $d < |D|$, $D_d = b$ and
$\fn{fresh}_T(k, b)$.  The range test is written as a nested conditional rather
than a conjunction so that $\fn{nth}\,d\,D$ is never mentioned when $d$ is out
of range.  The application rule instantiates it at $k = 2$, which is the sense
in which \BGA{} definitions are binary: the body may mention $v_0$ and $v_1$
and nothing else.

\subsection{Semantics}
\label{sec:fuel}

Two evaluators appear in the file.  The direct one, fixed by
\isa{bga\_semantics}, is the natural reading of the \BGA{} big-step semantics,
and writing it down at all depends on \QGA{} tolerating divergent definitions.

\thy{bga-eval}

\noindent
In mathematical notation, with $A$ an assignment,
\[
\fn{eval}(t, A) =
\begin{cases}
\fn{nth}\,j\,A                     & t = v_j\\
0                                  & t = 0\\
\Suc{\fn{eval}(a, A)}              & t = \Suc{a}\\
\Pre{\fn{eval}(a, A)}              & t = \Pre{a}\\
\fn{eval}(a, A) \text{ or } \fn{eval}(b, A)
                                   & t = (c\,0?\,a:b), \text{ on } \fn{eval}(c, A) = 0\\
\fn{eval}(D_i,\; [\fn{eval}(a, A),\, \fn{eval}(b, A)])
                                   & t = d_i(a, b)
\end{cases}
\]
The application case in the sources is wrapped in two guards of the form
\isa{eval a A = eval a A}.  Those are habeas quid tests: the arguments are
evaluated first, and the body is entered only if both terminate, which makes
the encoded semantics call by value.

The consistency argument runs through the step-indexed variant instead, fixed
by \isa{bga\_fuel\_semantics}.  Its result is offset by one, so that
$\fn{eval\_fuel}(k, t, A) = 0$ means ``no result within budget $k$'' and
$\fn{eval\_fuel}(k, t, A) = 1 + v$ means ``the value is $v$''.  The offset is
what lets a single \isa{num} carry both outcomes, at the cost of an explicit
$\Pre{\cdot}$ on every recursive result.

\thy{bga-eval-fuel}
\thy{bga-evals}

\noindent
\isa{evals t A r} says that $t$ evaluates to $r$ under $A$ for some budget:
\[
\fn{evals}(t, A, r) \;\;\equiv\;\; \exists k.\; \fn{eval\_fuel}(k, t, A) = \Suc{r}.
\]

Satisfaction comes in the same two forms.  For the direct evaluator,
\isa{bga\_semantics} defines

\thy{bga-sat}

\noindent
that is, $\fn{sat}(\varphi, A)$ compares $\fn{eval}$ of the two sides, with the
comparison chosen by the formula tag, and $\fn{sat\_hyp}$ quantifies over the
members of the list.  For the step-indexed evaluator, satisfaction of an
equation also has to require both sides to converge:

\thy{bga-sat-fuel}
\thy{bga-sat-hyp-fuel}

\[
\fn{sat\_fuel}(\varphi, A) \;\;\equiv\;\;
\exists x, y.\;\; \fn{evals}(a, A, x) \wedge \fn{evals}(b, A, y) \wedge
\begin{cases} x = y & \text{tag } \varphi = \fn{F\_EQ}\\
              x \neq y & \text{otherwise}\end{cases}
\]
where $a$ and $b$ are the two sides of $\varphi$.  A formula whose sides
diverge is unsatisfied under every assignment, which is the property the
consistency argument ultimately rests on.

\subsection{\isa{bga\_subst} and \isa{asn\_put}}

\thy{bga-subst}
\thy{bga-asn-put}

\begin{center}
\begin{tabular}{@{}ll@{}}
\toprule
\isa{subst\_T :: tm ⇒ tm ⇒ tm ⇒ tm} & $t[v_j \mapsto v]$, replace variable $j$ inside a term\\
\isa{subst\_F :: fm ⇒ tm ⇒ tm ⇒ fm} & the same inside a formula, side by side\\
\isa{subst\_body :: tm ⇒ tm ⇒ tm ⇒ tm} & $b[v_0 \mapsto x,\, v_1 \mapsto y]$, simultaneous\\
\isa{asn\_put :: asn ⇒ num ⇒ val ⇒ asn} & $A[i \mapsto v]$, update an assignment\\
\bottomrule
\end{tabular}
\end{center}

\noindent
All three substitutions recurse over the tag and rebuild with \isa{pack\_T}.
Two details are worth noting.  In the \isa{T\_ZERO} case \isa{subst\_T} returns
\isa{pack\_T T\_ZERO 0} rather than the original code, which normalizes the
ignored payload of a zero.  In the \isa{T\_APP} case the definition index
\isa{cpx (load\_T t)} is carried through untouched, since it is not a term.
\isa{subst\_body} is the two-variable version used by the application rule, and
it exists separately because that rule substitutes both arguments at once.
\isa{asn\_put A i v} walks down $i$ entries and replaces the one it finds,
padding with zeros when $i$ runs past the end of the list, so that an
assignment is a total map from indices to values.

\subsection{The proof rules of \BGA}
\label{sec:rules}

The rules are given below in sequent form, and each is realized as one branch
of the encoded checker.  Here $\Gamma$ is a list of encoded formulas,
$\varphi$ an encoded formula, and $\varphi[a]$ a formula with occurrences of
$a$ marked, which is realized by the template search of
\S\ref{sec:subst-rule}.  The turnstile is the \BGA{} sequent arrow, written
\isa{⊩} in the sources.

\begin{rules}
\iaxiom{\Gamma, \varphi \seq \varphi}{hyp} &
\iaxiom{\Gamma \seq 0 = 0}{0I}\\[2.4ex]
\irule{\Gamma \seq a = b}{\Gamma \seq b = a}{=S} &
\irule{\Gamma \seq a = b \sep \Gamma \seq \varphi[a]}{\Gamma \seq \varphi[b]}{=E}\\[2.4ex]
\irule{\Gamma \seq a = b}{\Gamma \seq \Suc{a} = \Suc{b}}{S=I} &
\irule{\Gamma \seq \Suc{a} = \Suc{b}}{\Gamma \seq a = b}{S=E}\\[2.4ex]
\irule{\Gamma \seq a \N}{\Gamma \seq \Pre{\Suc{a}} = a}{PI} &
\irule{\Gamma \seq a \N}{\Gamma \seq \Suc{a} \neq 0}{S{\neq}0I}\\[2.4ex]
\irule{\Gamma \seq c = 0 \sep \Gamma \seq a \N}{\Gamma \seq (c\,0?\,a:b) = a}{?I1} &
\irule{\Gamma \seq c \neq 0 \sep \Gamma \seq b \N}{\Gamma \seq (c\,0?\,a:b) = b}{?I2}\\[2.4ex]
\irule{\Gamma \seq a \neq b}{\Gamma \seq b \neq a}{{\neq}S} &
\irule{\Gamma \seq a \neq b}{\Gamma \seq \Suc{a} \neq \Suc{b}}{{\neq}SI}\\[2.4ex]
\irule{\Gamma \seq \Suc{a} \neq \Suc{b}}{\Gamma \seq a \neq b}{S{\neq}E} &
\irule{\Gamma \seq a \sep \Gamma, a \seq c}{\Gamma \seq c}{cut}\\[2.4ex]
\irule{\Gamma' \seq \varphi \sep \Gamma' \subseteq \Gamma}{\Gamma \seq \varphi}{wk} &
\irule{\begin{array}{c}
        \Gamma \seq a \N \sep \Gamma \seq \varphi[0]\\
        \Gamma, v_i \N, \varphi[v_i] \seq \varphi[\Suc{v_i}]
       \end{array}}{\Gamma \seq \varphi[a]}{Ind}\\[4.4ex]
\multicolumn{2}{c}{%
\irule{\Gamma \seq a_0 \N \sep \Gamma \seq a_1 \N \sep
       \Gamma \seq \varphi[\,D_d\langle a_0, a_1\rangle\,]}
      {\Gamma \seq \varphi[\,d(a_0, a_1)\,]}{app}}
\end{rules}

\noindent
Termination is not a separate judgment form in \BGA.  $a \N$ abbreviates the
formula $a = a$.  Almost every rule is an introduction rule, and the only
elimination is $=$E.

The checker never applies a rule forwards.  It is given a judgment $J$ and the
tail $\rho$ of the proof, and asks whether $J$ could have been produced by some
rule from judgments occurring in $\rho$.  Every rule with a premise therefore
becomes a search through $\rho$, and the rules whose premises mention a
$\varphi[\cdot]$ become a search through $\rho$ combined with a search for a
template.

\subsection{\isa{bga\_subst\_rule}}
\label{sec:subst-rule}

\begin{center}
\begin{tabular}{@{}ll@{}}
\toprule
\isa{rep\_vars\_T :: tm ⇒ tm ⇒ tm} & replaces every leaf of a term by a given variable\\
\isa{rep\_vars\_F :: fm ⇒ tm ⇒ fm} & the same on both sides of a formula\\
\isa{check\_template} & searches for a template below a bound\\
\isa{find\_phi} & searches $\rho$ for the premise $\Gamma \seq \varphi[a]$\\
\isa{find\_eq} & searches $\rho$ for the equation $\Gamma \seq a = b$\\
\isa{check\_subst :: jdg ⇒ pf ⇒ o} & the rule $=$E as a whole\\
\bottomrule
\end{tabular}
\end{center}

\noindent
The substitution rule is the one place where the checker has to invent
something.  To justify $\Gamma \seq \varphi[b]$ it must produce three things: a
premise $\Gamma \seq \varphi[a]$ somewhere in $\rho$, an equation
$\Gamma \seq a = b$ somewhere in $\rho$, and a \emph{template} $p$, a formula
with a designated variable $v_i$, such that
\[
p[v_i \mapsto a] = \varphi[a] \qquad\text{and}\qquad p[v_i \mapsto b] = \varphi[b].
\]
The template records which occurrences of $a$ are being rewritten.  Nothing in
the conclusion says what it is, so it is searched for.

\paragraph{Choosing the variable.}
The designated variable must not already occur in the judgment, or substituting
into the template would disturb something else.  The index chosen is $J + 1$,
where $J$ is the code of the whole judgment.  This is fresh because a variable
$v_j$ is encoded as $\fn{pack}_T(\fn{T\_VAR}, j)$, which is a pair with second
component $j$, and pairing and packing are monotone in that component, so
$j \leq J$ for every variable occurring anywhere in $J$.  Monotonicity of
\isa{pack} in \S\ref{sec:encoder-interface} is assumed for this argument.

\paragraph{Bounding the search.}
Templates are enumerated downwards from
$\fn{rep\_vars}_F(\varphi[b], J+1)$, the formula obtained by replacing every
leaf of the conclusion by $v_{J+1}$:
\[
\fn{rep\_vars}_T(t, i) =
\begin{cases}
v_i                                        & t = v_j \text{ or } t = 0\\
\Suc{\fn{rep\_vars}_T(a, i)}               & t = \Suc{a}\\
\Pre{\fn{rep\_vars}_T(a, i)}               & t = \Pre{a}\\
(\fn{rep\_vars}_T(c, i)\,0?\,\fn{rep\_vars}_T(a, i) : \fn{rep\_vars}_T(b, i))
                                           & t = (c\,0?\,a:b)\\
d_j(\fn{rep\_vars}_T(a, i), \fn{rep\_vars}_T(b, i))
                                           & t = d_j(a, b)
\end{cases}
\]
This is the template that abstracts \emph{every} leaf at once.  Its code
bounds the code of every other candidate with the same constructor skeleton,
since such a candidate differs from it only by leaving some leaves
unabstracted, and a leaf is numerically smaller than $v_{J+1}$.
\isa{check\_template} then counts $p$ down from that bound to zero and tests
each value:
\[
\fn{check\_template}(f, \varphi, a, b, p, i) \;\;\Longleftrightarrow\;\;
\exists q \leq p.\;\; q[v_i \mapsto a] = \varphi \;\wedge\; q[v_i \mapsto b] = f.
\]

\thy{bga-check-template}
\thy{bga-rep-vars-t}

\begin{remarkx}[The bound matters for completeness only]
Soundness of the checker does not depend on the bound being correct.
\isa{check\_template\_sound\_fuel} assumes that the search reported success and
extracts the witness $q$ it found; a bound that was too small would make the
checker reject proofs it ought to accept, which is a completeness failure, and
would leave soundness untouched.  Nothing in this development proves
completeness.
\end{remarkx}

\paragraph{Searching the proof.}
With the template search in place, the two remaining searches walk the proof
list.  In mathematical notation, writing $K \in \rho$ for membership in the
proof list,
\begin{align*}
\fn{find\_phi}(J, a, b, \pi) &\;\Longleftrightarrow\;
  \exists K \in \pi.\;\; \fn{hyp}(K) = \fn{hyp}(J) \;\wedge\;
  \fn{check\_template}(\fn{conc}(J), \fn{conc}(K), a, b, \beta_J, J+1),\\
\fn{find\_eq}(J, \rho, \pi) &\;\Longleftrightarrow\;
  \exists K \in \pi.\;\; \fn{hyp}(K) = \fn{hyp}(J) \;\wedge\;
  \fn{conc}(K) \text{ is } (a = b) \;\wedge\; \fn{find\_phi}(J, a, b, \rho),\\
\fn{check\_subst}(J, \rho) &\;\Longleftrightarrow\; \fn{find\_eq}(J, \rho, \rho),
\end{align*}
where $\beta_J = \fn{rep\_vars}_F(\fn{conc}(J), J+1)$ is the template bound.
Both searches carry two proof arguments.  The one named \isa{ptr} is the
part still to be scanned and shrinks on each step, and the one named
\isa{rest} is the whole tail of the proof and stays fixed, because the
premise being looked for may sit anywhere in it.  The hypothesis test
$\fn{hyp}(K) = \fn{hyp}(J)$ is what keeps the rule from silently changing the
context.

\thy{bga-find-phi}
\thy{bga-find-eq}

\subsection{\isa{bga\_eq\_rule}}

\begin{center}
\begin{tabular}{@{}ll@{}}
\toprule
\isa{check\_eq\_rules :: hyp ⇒ tm ⇒ tm ⇒ tmtag ⇒ tmtag ⇒ pf ⇒ o}
  & the seven rules concluding an equation\\
\isa{check\_neq\_rules :: hyp ⇒ tm ⇒ tm ⇒ tmtag ⇒ tmtag ⇒ pf ⇒ o}
  & the four rules concluding a disequation\\
\bottomrule
\end{tabular}
\end{center}

\noindent
Both take the two sides and their tags as separate arguments, already projected
out of the conclusion by \isa{valid\_step}, so that the branches can test tags
without recomputing them.

\thy{bga-eq-rules}

\noindent
The seven branches are, in order, \textsc{0I}, \textsc{=S}, \textsc{S=I},
\textsc{S=E}, \textsc{PI}, \textsc{?I2} and \textsc{?I1}.  Each is a membership
test on \isa{rest}, so a judgment is justified only by judgments occurring
later in the list, which is what makes \isa{check\_list} a well-founded
recursion.  Reading the fifth branch as an example: the conclusion is
$\Gamma \seq \Pre{u} = b$, the branch checks that $u$ is $\Suc{b}$ by
inspecting tags and payloads, and it requires $\Gamma \seq b = b$ to appear in
the proof, which is the habeas quid premise $b \N$ of \textsc{PI}.

\thy{bga-neq-rules}

\noindent
and dually \textsc{$\neq$S}, \textsc{S$\neq$0I}, \textsc{$\neq$SI} and
\textsc{S$\neq$E}.

\subsection{\isa{bga\_ind\_rule}}
\label{sec:ind-rule}

\begin{center}
\begin{tabular}{@{}ll@{}}
\toprule
\isa{check\_ind\_template} & tests one candidate template for the induction rule\\
\isa{find\_ind\_base} & searches $\rho$ for the base case\\
\isa{check\_ind :: jdg ⇒ pf ⇒ o} & the rule \textsc{Ind} as a whole\\
\bottomrule
\end{tabular}
\end{center}

\noindent
Induction needs the same kind of template, and three premises rather than one.
To justify $\Gamma \seq \varphi[a]$ it looks for a template $q$ with
\[
q[v_i \mapsto a] = \varphi[a], \qquad
q[v_i \mapsto 0] \text{ present in } \rho, \qquad
(\,v_i = v_i,\; q,\; \Gamma \seq q[v_i \mapsto \Suc{v_i}]\,) \text{ present in } \rho,
\]
that is,
\begin{align*}
\fn{check\_ind\_template}(f, \varphi, a, \Gamma, p, i, \rho)
  &\;\Longleftrightarrow\; \exists q \leq p.\;\;
   q[v_i \mapsto a] = f \;\wedge\; q[v_i \mapsto 0] = \varphi\\
  &\qquad\wedge\;\;
   \bigl(v_i = v_i,\, q,\, \Gamma \seq q[v_i \mapsto \Suc{v_i}]\bigr) \in \rho.
\end{align*}

\thy{bga-ind-template}
\thy{bga-check-ind}

\noindent
\isa{check\_ind} guards the search with two conditions.  The habeas quid
premise $\Gamma \seq a \N$, encoded as the judgment $\Gamma \seq a = a$, must
be in the proof, and $\fn{fresh}_H(J+1, \Gamma)$ must hold, so that the
variable used to build the template does not occur in the hypotheses.  The
step case sought in the proof carries both the induction hypothesis $q$ and the
termination assumption $v_i = v_i$ on the induction variable, which is what the
semantic induction lemma of \S\ref{sec:proof} consumes.  \isa{find\_ind\_base}
scans $\rho$ for the base case in the same style as \isa{find\_phi}.

\subsection{\isa{bga\_struct\_rule}}

\begin{center}
\begin{tabular}{@{}ll@{}}
\toprule
\isa{find\_cut} & searches $\rho$ for the cut formula\\
\isa{check\_cut :: jdg ⇒ pf ⇒ o} & the rule \textsc{cut}\\
\isa{find\_struct} & searches $\rho$ for the weakened judgment\\
\isa{check\_struct :: jdg ⇒ pf ⇒ o} & the rule \textsc{wk}\\
\bottomrule
\end{tabular}
\end{center}

\thy{bga-struct}

\begin{align*}
\fn{check\_cut}(J, \rho) &\;\Longleftrightarrow\;
  \exists K \in \rho.\;\; \fn{hyp}(K) = \fn{hyp}(J) \;\wedge\;
  \bigl(\fn{conc}(K),\, \fn{hyp}(J) \seq \fn{conc}(J)\bigr) \in \rho,\\
\fn{check\_struct}(J, \rho) &\;\Longleftrightarrow\;
  \exists K \in \rho.\;\; \fn{conc}(K) = \fn{conc}(J) \;\wedge\;
  \fn{hyp}(K) \subseteq \fn{hyp}(J).
\end{align*}

\noindent
The hypothesis rule is not a separate function.  It is the first branch of
\isa{valid\_step}, the test \isa{mem (conc\_of J) (hyp\_of J)}.

\subsection{\isa{bga\_app\_rule}}

\begin{center}
\begin{tabular}{@{}ll@{}}
\toprule
\isa{app\_try} & tests one triple $(d, x, y)$\\
\isa{app\_y} & counts $y$ down\\
\isa{app\_x} & counts $x$ down\\
\isa{app\_d} & counts the definition index $d$ down\\
\isa{check\_app :: jdg ⇒ pf ⇒ o} & the rule \textsc{app} as a whole\\
\bottomrule
\end{tabular}
\end{center}

\thy{bga-app}

\noindent
Application unfolds a definition.  To justify $\Gamma \seq \varphi[d(x,y)]$ the
checker needs $\Gamma \seq x \N$, $\Gamma \seq y \N$ and a premise
$\Gamma \seq \varphi[\,D_d[v_0 \mapsto x, v_1 \mapsto y]\,]$:
\begin{align*}
\fn{app\_try}(J, d, x, y, \rho) \;\Longleftrightarrow\;\;
  & \fn{dfn\_is}(d, 2, D_d)
    \;\wedge\; (\fn{hyp}(J) \seq x = x) \in \rho
    \;\wedge\; (\fn{hyp}(J) \seq y = y) \in \rho\\
  &\wedge\; \fn{find\_phi}\bigl(J,\; D_d[v_0 \mapsto x, v_1 \mapsto y],\; d(x,y),\; \rho\bigr).
\end{align*}
Note that \isa{find\_phi} is reused here with the two sides of the rewrite
being the unfolded body and the application, so unfolding a definition is
treated as one more instance of rewriting inside a template.

None of $d$, $x$ and $y$ is determined by the conclusion, so all three are
enumerated:
\[
\fn{check\_app}(J, \rho) \;\Longleftrightarrow\;
\exists d < |D|.\;\; \exists x \leq \fn{conc}(J).\;\; \exists y \leq \fn{conc}(J).\;\;
\fn{app\_try}(J, d, x, y, \rho).
\]
The three nested countdowns \isa{app\_d}, \isa{app\_x} and \isa{app\_y} are
that triple search written as recursion, since \QGA{} has no bounded
quantifier.  The bound $\fn{conc}(J)$ on the arguments is the same kind of
bound as in \S\ref{sec:subst-rule}: the code of a subterm is at most the code
of the formula containing it, and as before it is completeness rather than
soundness that depends on it.

\subsection{\isa{bga\_proof\_check}}

\begin{center}
\begin{tabular}{@{}ll@{}}
\toprule
\isa{valid\_step :: jdg ⇒ pf ⇒ o} & ``$J$ follows by one rule from judgments in $\rho$''\\
\isa{check\_list :: pf ⇒ o} & every entry follows from the entries after it\\
\isa{is\_valid\_proof :: pf ⇒ jdg ⇒ o} & ``$\pi$ is a valid proof whose head is $J$''\\
\bottomrule
\end{tabular}
\end{center}

\thy{bga-valid-step}

\noindent
\isa{valid\_step} is the disjunction of the branches, tried in order:
membership for \textsc{hyp}, then cut, substitution, induction, application and
weakening, and finally the equality or disequality rules according to the tag
of the conclusion.  Written as a formula,
\begin{align*}
\fn{valid\_step}(J, \rho) \;\Longleftrightarrow\;\;
& \fn{conc}(J) \in \fn{hyp}(J)
  \;\vee\; \fn{check\_cut}(J, \rho)
  \;\vee\; \fn{check\_subst}(J, \rho)\\
&\vee\; \fn{check\_ind}(J, \rho)
  \;\vee\; \fn{check\_app}(J, \rho)
  \;\vee\; \fn{check\_struct}(J, \rho)\\
&\vee\; \bigl(\text{tag } \fn{conc}(J) = \fn{F\_EQ}
   \;?\; \fn{check\_eq\_rules}(\dots) : \fn{check\_neq\_rules}(\dots)\bigr).
\end{align*}
The order is not significant for soundness.  It is a sequence of nested
conditionals rather than a disjunction because each branch has to be evaluated
only when the earlier ones fail.

\thy{bga-check-list}
\thy{bga-is-valid-proof}

\noindent
A proof is valid when every entry follows from the entries after it, and it
proves $J$ when additionally its head is $J$:
\[
\fn{check\_list}(\pi) \;\Longleftrightarrow\;
\pi = \fn{Nil} \;\vee\;
\bigl(\fn{valid\_step}(\fn{hd}\,\pi,\, \fn{tl}\,\pi) \wedge \fn{check\_list}(\fn{tl}\,\pi)\bigr),
\]
\[
\fn{is\_valid\_proof}(\pi, J) \;\Longleftrightarrow\;
\pi \neq \fn{Nil} \;\wedge\; \fn{hd}\,\pi = J \;\wedge\; \fn{check\_list}(\pi).
\]
Reading the list head first, with each entry justified by its tail, is what
makes the recursion terminate and what the soundness induction of
\S\ref{sec:list} follows.
